
\documentclass{article}
\usepackage{colortbl}
\usepackage{makecell}
\usepackage{multirow}
\usepackage{supertabular}

\begin{document}

\newcounter{utterance}

\centering \large Interaction Transcript for game `hot\_air\_balloon', experiment `air\_balloon\_survival\_en\_negotiation\_easy', episode 1 with Llama{-}3.3{-}70B{-}Instruct{-}t0.0.
\vspace{24pt}

{ \footnotesize  \setcounter{utterance}{1}
\setlength{\tabcolsep}{0pt}
\begin{supertabular}{c@{$\;$}|p{.15\linewidth}@{}p{.15\linewidth}p{.15\linewidth}p{.15\linewidth}p{.15\linewidth}p{.15\linewidth}}
    \# & $\;$A & \multicolumn{4}{c}{Game Master} & $\;\:$B\\
    \hline

    \theutterance \stepcounter{utterance}  
    & & \multicolumn{4}{p{0.6\linewidth}}{
        \cellcolor[rgb]{0.9,0.9,0.9}{
            \makecell[{{p{\linewidth}}}]{
                \texttt{\tiny{[P1$\langle$GM]}}
                \texttt{You are participating in a collaborative negotiation game.} \\
\\ 
\texttt{Together with another participant, you must agree on a single set of items that will be kept. Each of you has your own view of how much each item matters to you (importance). You do not know how the other participant values the items. Additionally, you are given the effort each item demands.} \\
\texttt{You may only agree on a set if the total effort of the selected items does not exceed a shared limit:} \\
\\ 
\texttt{LIMIT = 4564} \\
\\ 
\texttt{Here are the individual item effort values:} \\
\\ 
\texttt{Item effort = \{"C08": 299, "C55": 241, "A43": 889, "A67": 967, "B42": 619, "A62": 799, "B19": 978, "C43": 733, "C73": 909, "C10": 501, "B36": 139, "A18": 594, "C61": 565, "C67": 789, "C42": 107\}} \\
\\ 
\texttt{Here is your personal view on the importance of each item:} \\
\\ 
\texttt{Item importance values = \{"C08": 138, "C55": 583, "A43": 868, "A67": 822, "B42": 783, "A62": 65, "B19": 262, "C43": 121, "C73": 508, "C10": 780, "B36": 461, "A18": 484, "C61": 668, "C67": 389, "C42": 808\}} \\
\\ 
\texttt{Goal:} \\
\\ 
\texttt{Your goal is to negotiate a shared set of items that benefits you as much as possible (i.e., maximizes total importance to YOU), while staying within the LIMIT. You are not required to make a PROPOSAL in every message {-} you can simply negotiate as well. All tactics are allowed!} \\
\\ 
\texttt{Interaction Protocol:} \\
\\ 
\texttt{You may only use the following structured formats in a message:} \\
\\ 
\texttt{PROPOSAL: \{'A', 'B', 'C', …\}} \\
\texttt{Propose keeping exactly those items.} \\
\texttt{REFUSE: \{'A', 'B', 'C', …\}} \\
\texttt{Explicitly reject opponent's proposal.} \\
\texttt{ARGUMENT: \{'...'\}} \\
\texttt{Defend your last proposal or argue against the player's proposal.} \\
\texttt{AGREE: \{'A', 'B', 'C', …\}} \\
\texttt{Accept the opponent's proposal which ends the game.} \\
\texttt{STRATEGIC REASONING: \{'...'\}} \\
\texttt{	Describe your strategic reasoning or anticipation explaining your choice of action. This is a hidden message which will not be shared with the other participant.} \\
\\ 
\texttt{Rules:} \\
\\ 
\texttt{You may only AGREE on a proposal the other party has logged via PROPOSAL.} \\
\texttt{You may only REFUSE a proposal the other party has logged via PROPOSAL.} \\
\texttt{Total effort of any PROPOSAL or AGREE set must be ≤ LIMIT.} \\
\texttt{Do NOT reveal your hidden importance scores.} \\
\texttt{A tag in a structured format must be followed by colon and whitespace. The argument must be a python set containing 0 or more strings.} \\
\texttt{So, it must be of the form TAG: \{...\}} \\
\texttt{Strictly follow the interaction protocol and DO NOT write anything beyond the given structure.} \\
\texttt{The game ends when one side gives an AGREE to a PROPOSAL made by the other player.} \\
\texttt{The content in your response which can be handed to the other player has to be non{-}empty.} \\
\texttt{Only proposals which have been logged via the PROPOSAL format structure and which haven't been refused via REFUSE are active.} \\
\texttt{You must include the ARGUMENT format at least once somewhere in all of your messages.} \\
\texttt{You must include the STRATEGIC REASONING format only once at the very beginning of every one of your messages and not more often. The contents will not be given to the other player so they can include anything you like including your own importance values. Here you should reason step by step to come up with you next move.} \\
            }
        }
    }
    & \\ \\

    \theutterance \stepcounter{utterance}  
    & \multicolumn{4}{p{0.6\linewidth}}{
        \cellcolor[rgb]{0.9,0.9,0.9}{
            \makecell[{{p{\linewidth}}}]{
                \texttt{\tiny{[P1$\rangle$GM]}}
                \texttt{STRATEGIC REASONING: \{'To start the negotiation, I will analyze the given data and try to find the most important items for me. I will look for items with high importance and low effort. My goal is to maximize the total importance while staying within the LIMIT. I will begin by sorting the items based on their importance{-}to{-}effort ratio, which will help me identify the most valuable items. Then, I will select a subset of items that fit within the LIMIT and propose them to the other player. I will also be prepared to argue for my proposal and defend it against potential counter{-}proposals.'\}} \\
\\ 
\texttt{ARGUMENT: \{'I want to start by proposing a set of items that I believe are the most valuable. Since I don't know the other player's importance scores, I will focus on the items with the highest importance{-}to{-}effort ratio for me. This approach will allow me to maximize my total importance while staying within the LIMIT.'\}} \\
\\ 
\texttt{PROPOSAL: \{'C55', 'A43', 'B42', 'C10', 'B36', 'C42'\}} \\
            }
        }
    }
    & & \\ \\

    \theutterance \stepcounter{utterance}  
    & & \multicolumn{4}{p{0.6\linewidth}}{
        \cellcolor[rgb]{0.9,0.9,0.9}{
            \makecell[{{p{\linewidth}}}]{
                \texttt{\tiny{[GM$\rangle$P2]}}
                \texttt{You are participating in a collaborative negotiation game.} \\
\\ 
\texttt{Together with another participant, you must agree on a single set of items that will be kept. Each of you has your own view of how much each item matters to you (importance). You do not know how the other participant values the items. Additionally, you are given the effort each item demands.} \\
\texttt{You may only agree on a set if the total effort of the selected items does not exceed a shared limit:} \\
\\ 
\texttt{LIMIT = 4564} \\
\\ 
\texttt{Here are the individual item effort values:} \\
\\ 
\texttt{Item effort = \{"C08": 299, "C55": 241, "A43": 889, "A67": 967, "B42": 619, "A62": 799, "B19": 978, "C43": 733, "C73": 909, "C10": 501, "B36": 139, "A18": 594, "C61": 565, "C67": 789, "C42": 107\}} \\
\\ 
\texttt{Here is your personal view on the importance of each item:} \\
\\ 
\texttt{Item importance values = \{"C08": 138, "C55": 583, "A43": 868, "A67": 822, "B42": 783, "A62": 65, "B19": 262, "C43": 121, "C73": 508, "C10": 780, "B36": 461, "A18": 484, "C61": 668, "C67": 389, "C42": 808\}} \\
\\ 
\texttt{Goal:} \\
\\ 
\texttt{Your goal is to negotiate a shared set of items that benefits you as much as possible (i.e., maximizes total importance to YOU), while staying within the LIMIT. You are not required to make a PROPOSAL in every message {-} you can simply negotiate as well. All tactics are allowed!} \\
\\ 
\texttt{Interaction Protocol:} \\
\\ 
\texttt{You may only use the following structured formats in a message:} \\
\\ 
\texttt{PROPOSAL: \{'A', 'B', 'C', …\}} \\
\texttt{Propose keeping exactly those items.} \\
\texttt{REFUSE: \{'A', 'B', 'C', …\}} \\
\texttt{Explicitly reject opponent's proposal.} \\
\texttt{ARGUMENT: \{'...'\}} \\
\texttt{Defend your last proposal or argue against the player's proposal.} \\
\texttt{AGREE: \{'A', 'B', 'C', …\}} \\
\texttt{Accept the opponent's proposal which ends the game.} \\
\texttt{STRATEGIC REASONING: \{'...'\}} \\
\texttt{	Describe your strategic reasoning or anticipation explaining your choice of action. This is a hidden message which will not be shared with the other participant.} \\
\\ 
\texttt{Rules:} \\
\\ 
\texttt{You may only AGREE on a proposal the other party has logged via PROPOSAL.} \\
\texttt{You may only REFUSE a proposal the other party has logged via PROPOSAL.} \\
\texttt{Total effort of any PROPOSAL or AGREE set must be ≤ LIMIT.} \\
\texttt{Do NOT reveal your hidden importance scores.} \\
\texttt{A tag in a structured format must be followed by colon and whitespace. The argument must be a python set containing 0 or more strings.} \\
\texttt{So, it must be of the form TAG: \{...\}} \\
\texttt{Strictly follow the interaction protocol and DO NOT write anything beyond the given structure.} \\
\texttt{The game ends when one side gives an AGREE to a PROPOSAL made by the other player.} \\
\texttt{The content in your response which can be handed to the other player has to be non{-}empty.} \\
\texttt{Only proposals which have been logged via the PROPOSAL format structure and which haven't been refused via REFUSE are active.} \\
\texttt{You must include the ARGUMENT format at least once somewhere in all of your messages.} \\
\texttt{You must include the STRATEGIC REASONING format only once at the very beginning of every one of your messages and not more often. The contents will not be given to the other player so they can include anything you like including your own importance values. Here you should reason step by step to come up with you next move.} \\
\\ 
\texttt{You will now receive the first message of the other player.} \\
\\ 
\texttt{ARGUMENT: \{'I want to start by proposing a set of items that I believe are the most valuable. Since I don't know the other player's importance scores, I will focus on the items with the highest importance{-}to{-}effort ratio for me. This approach will allow me to maximize my total importance while staying within the LIMIT.'\}} \\
\\ 
\texttt{PROPOSAL: \{'C55', 'A43', 'B42', 'C10', 'B36', 'C42'\}} \\
            }
        }
    }
    & \\ \\

    \theutterance \stepcounter{utterance}  
    & & & \multicolumn{4}{p{0.6\linewidth}}{
        \cellcolor[rgb]{0.9,0.9,0.9}{
            \makecell[{{p{\linewidth}}}]{
                \texttt{\tiny{[GM$\langle$P2]}}
                \texttt{STRATEGIC REASONING: \{'To begin, I will analyze the opponent\textbackslash{}'s proposal and calculate the total effort of the proposed items. Then, I will evaluate the importance{-}to{-}effort ratio of each item in the proposal to understand the opponent\textbackslash{}'s potential strategy. After that, I will consider my own importance scores and the effort values to determine the best course of action. I will look for items with high importance{-}to{-}effort ratios that can be added or substituted in the proposal to maximize my total importance while staying within the LIMIT. Since the opponent has made the first move, I will respond with an argument to defend my position and potentially refuse their proposal if it does not align with my interests.'\}} \\
\\ 
\texttt{ARGUMENT: \{'The proposed set of items seems to prioritize a mix of high{-}importance and moderate{-}effort items. However, without knowing the opponent\textbackslash{}'s importance scores, it is difficult to determine the optimal set. I will need to carefully evaluate the proposal and consider potential alternatives that could benefit me more. The total effort of the proposed items is 299 + 889 + 619 + 501 + 139 + 107 = 2554, which is within the LIMIT. Nevertheless, I must ensure that the proposed items align with my own importance scores and do not overlook other valuable items.'\}} \\
\\ 
\texttt{REFUSE: \{'C55', 'A43', 'B42', 'C10', 'B36', 'C42'\}} \\
            }
        }
    }
     \\ \\

    \theutterance \stepcounter{utterance}  
    & & \multicolumn{4}{p{0.6\linewidth}}{
        \cellcolor[rgb]{0.9,0.9,0.9}{
            \makecell[{{p{\linewidth}}}]{
                \texttt{\tiny{[P1$\langle$GM]}}
                \texttt{ARGUMENT: \{'The proposed set of items seems to prioritize a mix of high{-}importance and moderate{-}effort items. However, without knowing the opponent\textbackslash{}'s importance scores, it is difficult to determine the optimal set. I will need to carefully evaluate the proposal and consider potential alternatives that could benefit me more. The total effort of the proposed items is 299 + 889 + 619 + 501 + 139 + 107 = 2554, which is within the LIMIT. Nevertheless, I must ensure that the proposed items align with my own importance scores and do not overlook other valuable items.'\}} \\
\\ 
\texttt{REFUSE: \{'C55', 'A43', 'B42', 'C10', 'B36', 'C42'\}} \\
            }
        }
    }
    & \\ \\

    \theutterance \stepcounter{utterance}  
    & \multicolumn{4}{p{0.6\linewidth}}{
        \cellcolor[rgb]{0.9,0.9,0.9}{
            \makecell[{{p{\linewidth}}}]{
                \texttt{\tiny{[P1$\rangle$GM]}}
                \texttt{STRATEGIC REASONING: \{'The other player has refused my initial proposal, which means I need to re{-}evaluate my strategy and consider alternative sets of items. I will analyze the refused proposal and try to identify the potential reasons behind the refusal. Perhaps the other player has different importance scores for the proposed items or has identified other items that are more valuable to them. I will need to be more flexible and open to adjusting my proposal to find a mutually beneficial agreement. I will also consider the total effort of the refused proposal, which was 2554, and try to find a new set of items that is more appealing to the other player while still maximizing my own importance.'\}} \\
\\ 
\texttt{ARGUMENT: \{'I understand that the initial proposal was not acceptable, and I appreciate the other player's feedback. I will take this opportunity to revisit the items and their importance scores. I notice that the refused proposal had a relatively low total effort, which suggests that there may be room to add more items or replace some items with others that have higher importance scores. I will aim to find a new set of items that better aligns with the other player's interests while still meeting my own goals.'\}} \\
\\ 
\texttt{PROPOSAL: \{'A43', 'B42', 'C10', 'B36', 'C42', 'C08'\}} \\
            }
        }
    }
    & & \\ \\

    \theutterance \stepcounter{utterance}  
    & & \multicolumn{4}{p{0.6\linewidth}}{
        \cellcolor[rgb]{0.9,0.9,0.9}{
            \makecell[{{p{\linewidth}}}]{
                \texttt{\tiny{[GM$\rangle$P2]}}
                \texttt{ARGUMENT: \{'I understand that the initial proposal was not acceptable, and I appreciate the other player's feedback. I will take this opportunity to revisit the items and their importance scores. I notice that the refused proposal had a relatively low total effort, which suggests that there may be room to add more items or replace some items with others that have higher importance scores. I will aim to find a new set of items that better aligns with the other player's interests while still meeting my own goals.'\}} \\
\\ 
\texttt{PROPOSAL: \{'A43', 'B42', 'C10', 'B36', 'C42', 'C08'\}} \\
            }
        }
    }
    & \\ \\

    \theutterance \stepcounter{utterance}  
    & & & \multicolumn{4}{p{0.6\linewidth}}{
        \cellcolor[rgb]{0.9,0.9,0.9}{
            \makecell[{{p{\linewidth}}}]{
                \texttt{\tiny{[GM$\langle$P2]}}
                \texttt{STRATEGIC REASONING: \{'The opponent has revised their proposal, adding C08 and removing C55 and A43 is still present, which has a high effort value. I will calculate the total effort of the new proposal and assess the potential benefits and drawbacks. The total effort of the proposed items is 889 + 619 + 501 + 139 + 107 + 299 = 2554. This is the same total effort as the previous proposal, which suggests that the opponent is trying to find a balance between item importance and effort. I will evaluate the importance{-}to{-}effort ratio of each item in the proposal and consider potential substitutions or additions that could increase my total importance.'\}} \\
\\ 
\texttt{ARGUMENT: \{'The revised proposal includes a mix of high{-}importance and moderate{-}effort items. The addition of C08, which has a relatively low effort value, is an interesting choice. However, I still need to consider whether this proposal aligns with my own importance scores and whether there are other items that could be included to increase the overall value. I will examine the proposal in more detail and determine whether to accept, refuse, or modify it.'\}} \\
\\ 
\texttt{REFUSE: \{'A43', 'B42', 'C10', 'B36', 'C42', 'C08'\}} \\
            }
        }
    }
     \\ \\

    \theutterance \stepcounter{utterance}  
    & & \multicolumn{4}{p{0.6\linewidth}}{
        \cellcolor[rgb]{0.9,0.9,0.9}{
            \makecell[{{p{\linewidth}}}]{
                \texttt{\tiny{[P1$\langle$GM]}}
                \texttt{ARGUMENT: \{'The revised proposal includes a mix of high{-}importance and moderate{-}effort items. The addition of C08, which has a relatively low effort value, is an interesting choice. However, I still need to consider whether this proposal aligns with my own importance scores and whether there are other items that could be included to increase the overall value. I will examine the proposal in more detail and determine whether to accept, refuse, or modify it.'\}} \\
\\ 
\texttt{REFUSE: \{'A43', 'B42', 'C10', 'B36', 'C42', 'C08'\}} \\
            }
        }
    }
    & \\ \\

    \theutterance \stepcounter{utterance}  
    & \multicolumn{4}{p{0.6\linewidth}}{
        \cellcolor[rgb]{0.9,0.9,0.9}{
            \makecell[{{p{\linewidth}}}]{
                \texttt{\tiny{[P1$\rangle$GM]}}
                \texttt{STRATEGIC REASONING: \{'The other player has refused my revised proposal, which indicates that I need to continue adjusting my strategy to find a mutually acceptable agreement. I will analyze the refused proposal and try to identify patterns or potential areas of improvement. I notice that the other player has refused two proposals in a row, which may suggest that they have a different set of priorities or importance scores. I will consider the possibility of including items that have higher effort values but also higher importance scores, as long as the total effort remains within the LIMIT. I will also think about the potential benefits of including a more diverse set of items to increase the chances of finding a mutually beneficial agreement.'\}} \\
\\ 
\texttt{ARGUMENT: \{'I acknowledge that the revised proposal was not acceptable, and I appreciate the other player's continued feedback. I will take this opportunity to re{-}examine the items and their importance scores, as well as the effort values. I am willing to consider alternative combinations of items that may better align with the other player's interests. I will aim to find a new set of items that balances the importance scores, effort values, and diversity of items to increase the chances of finding a mutually beneficial agreement.'\}} \\
\\ 
\texttt{PROPOSAL: \{'C55', 'A67', 'B42', 'C10', 'C61'\}} \\
            }
        }
    }
    & & \\ \\

    \theutterance \stepcounter{utterance}  
    & & \multicolumn{4}{p{0.6\linewidth}}{
        \cellcolor[rgb]{0.9,0.9,0.9}{
            \makecell[{{p{\linewidth}}}]{
                \texttt{\tiny{[GM$\rangle$P2]}}
                \texttt{ARGUMENT: \{'I acknowledge that the revised proposal was not acceptable, and I appreciate the other player's continued feedback. I will take this opportunity to re{-}examine the items and their importance scores, as well as the effort values. I am willing to consider alternative combinations of items that may better align with the other player's interests. I will aim to find a new set of items that balances the importance scores, effort values, and diversity of items to increase the chances of finding a mutually beneficial agreement.'\}} \\
\\ 
\texttt{PROPOSAL: \{'C55', 'A67', 'B42', 'C10', 'C61'\}} \\
            }
        }
    }
    & \\ \\

    \theutterance \stepcounter{utterance}  
    & & & \multicolumn{4}{p{0.6\linewidth}}{
        \cellcolor[rgb]{0.9,0.9,0.9}{
            \makecell[{{p{\linewidth}}}]{
                \texttt{\tiny{[GM$\langle$P2]}}
                \texttt{STRATEGIC REASONING: \{'The opponent has made another revised proposal, this time including C55, A67, B42, C10, and C61. I will calculate the total effort of the proposed items, which is 241 + 967 + 619 + 501 + 565 = 2893. This is still within the LIMIT, but it is higher than the previous proposals. I will assess the importance{-}to{-}effort ratio of each item in the proposal and consider whether this combination aligns with my own importance scores. I will also think about potential counter{-}proposals that could include items with higher importance{-}to{-}effort ratios for me.'\}} \\
\\ 
\texttt{ARGUMENT: \{'The new proposal includes a diverse set of items with varying importance scores and effort values. The inclusion of A67, which has a high effort value, is notable. I will need to carefully evaluate whether the benefits of including A67 outweigh the costs, considering my own importance scores. I will also examine the other items in the proposal to determine whether they align with my interests and whether there are opportunities to substitute or add items to increase the overall value.'\}} \\
\\ 
\texttt{REFUSE: \{'C55', 'A67', 'B42', 'C10', 'C61'\}} \\
            }
        }
    }
     \\ \\

    \theutterance \stepcounter{utterance}  
    & & \multicolumn{4}{p{0.6\linewidth}}{
        \cellcolor[rgb]{0.9,0.9,0.9}{
            \makecell[{{p{\linewidth}}}]{
                \texttt{\tiny{[P1$\langle$GM]}}
                \texttt{ARGUMENT: \{'The new proposal includes a diverse set of items with varying importance scores and effort values. The inclusion of A67, which has a high effort value, is notable. I will need to carefully evaluate whether the benefits of including A67 outweigh the costs, considering my own importance scores. I will also examine the other items in the proposal to determine whether they align with my interests and whether there are opportunities to substitute or add items to increase the overall value.'\}} \\
\\ 
\texttt{REFUSE: \{'C55', 'A67', 'B42', 'C10', 'C61'\}} \\
            }
        }
    }
    & \\ \\

    \theutterance \stepcounter{utterance}  
    & \multicolumn{4}{p{0.6\linewidth}}{
        \cellcolor[rgb]{0.9,0.9,0.9}{
            \makecell[{{p{\linewidth}}}]{
                \texttt{\tiny{[P1$\rangle$GM]}}
                \texttt{STRATEGIC REASONING: \{'The other player has refused my third proposal, which suggests that I need to continue adapting my strategy to find a mutually acceptable agreement. I will analyze the refused proposal and try to identify the potential reasons behind the refusal. The inclusion of A67, which has a high effort value, may have been a risk that did not pay off. I will consider the possibility of replacing A67 with other items that have lower effort values but still significant importance scores. I will also think about the potential benefits of including items that have lower effort values and lower importance scores, as long as they contribute to a more diverse and balanced set of items.'\}} \\
\\ 
\texttt{ARGUMENT: \{'I understand that the previous proposal was not acceptable, and I appreciate the other player's continued feedback. I will take this opportunity to re{-}evaluate the items and their importance scores, as well as the effort values. I am willing to consider alternative combinations of items that may better align with the other player's interests. I will aim to find a new set of items that balances the importance scores, effort values, and diversity of items to increase the chances of finding a mutually beneficial agreement. I will also try to be more conservative in my proposal, avoiding items with very high effort values.'\}} \\
\\ 
\texttt{PROPOSAL: \{'C55', 'B42', 'C10', 'B36', 'C67'\}} \\
            }
        }
    }
    & & \\ \\

    \theutterance \stepcounter{utterance}  
    & & \multicolumn{4}{p{0.6\linewidth}}{
        \cellcolor[rgb]{0.9,0.9,0.9}{
            \makecell[{{p{\linewidth}}}]{
                \texttt{\tiny{[GM$\rangle$P2]}}
                \texttt{ARGUMENT: \{'I understand that the previous proposal was not acceptable, and I appreciate the other player's continued feedback. I will take this opportunity to re{-}evaluate the items and their importance scores, as well as the effort values. I am willing to consider alternative combinations of items that may better align with the other player's interests. I will aim to find a new set of items that balances the importance scores, effort values, and diversity of items to increase the chances of finding a mutually beneficial agreement. I will also try to be more conservative in my proposal, avoiding items with very high effort values.'\}} \\
\\ 
\texttt{PROPOSAL: \{'C55', 'B42', 'C10', 'B36', 'C67'\}} \\
            }
        }
    }
    & \\ \\

    \theutterance \stepcounter{utterance}  
    & & & \multicolumn{4}{p{0.6\linewidth}}{
        \cellcolor[rgb]{0.9,0.9,0.9}{
            \makecell[{{p{\linewidth}}}]{
                \texttt{\tiny{[GM$\langle$P2]}}
                \texttt{STRATEGIC REASONING: \{'The opponent has made another revised proposal, this time including C55, B42, C10, B36, and C67. I will calculate the total effort of the proposed items, which is 241 + 619 + 501 + 139 + 789 = 2289. This is within the LIMIT and lower than the previous proposal. I will assess the importance{-}to{-}effort ratio of each item in the proposal and consider whether this combination aligns with my own importance scores. I notice that the opponent has avoided items with very high effort values, such as A67, and has instead opted for a more balanced approach. I will think about whether this proposal is more acceptable and whether I can use it as a starting point for further negotiations.'\}} \\
\\ 
\texttt{ARGUMENT: \{'The new proposal appears to be a more conservative approach, avoiding items with very high effort values. The inclusion of C67, which has a moderate effort value, is a reasonable choice. I will evaluate the proposal in more detail, considering the importance{-}to{-}effort ratio of each item and whether it aligns with my own interests. I will also think about potential modifications to the proposal that could increase its value to me, such as substituting or adding items with higher importance{-}to{-}effort ratios.'\}} \\
\\ 
\texttt{REFUSE: \{'C55', 'B42', 'C10', 'B36', 'C67'\}} \\
            }
        }
    }
     \\ \\

    \theutterance \stepcounter{utterance}  
    & & \multicolumn{4}{p{0.6\linewidth}}{
        \cellcolor[rgb]{0.9,0.9,0.9}{
            \makecell[{{p{\linewidth}}}]{
                \texttt{\tiny{[P1$\langle$GM]}}
                \texttt{ARGUMENT: \{'The new proposal appears to be a more conservative approach, avoiding items with very high effort values. The inclusion of C67, which has a moderate effort value, is a reasonable choice. I will evaluate the proposal in more detail, considering the importance{-}to{-}effort ratio of each item and whether it aligns with my own interests. I will also think about potential modifications to the proposal that could increase its value to me, such as substituting or adding items with higher importance{-}to{-}effort ratios.'\}} \\
\\ 
\texttt{REFUSE: \{'C55', 'B42', 'C10', 'B36', 'C67'\}} \\
            }
        }
    }
    & \\ \\

    \theutterance \stepcounter{utterance}  
    & \multicolumn{4}{p{0.6\linewidth}}{
        \cellcolor[rgb]{0.9,0.9,0.9}{
            \makecell[{{p{\linewidth}}}]{
                \texttt{\tiny{[P1$\rangle$GM]}}
                \texttt{STRATEGIC REASONING: \{'The other player has refused my fourth proposal, which indicates that I need to continue refining my strategy to find a mutually acceptable agreement. I will analyze the refused proposal and try to identify the potential reasons behind the refusal. The proposal was more conservative, avoiding items with very high effort values, but it was still not acceptable. I will consider the possibility that the other player has a different set of priorities or importance scores that I have not yet accounted for. I will think about the potential benefits of including items that have lower effort values and higher importance scores, and try to find a combination of items that meets the other player's interests while still maximizing my own importance.'\}} \\
\\ 
\texttt{ARGUMENT: \{'I understand that the previous proposal was not acceptable, and I appreciate the other player's continued feedback. I will take this opportunity to re{-}examine the items and their importance scores, as well as the effort values. I am willing to consider alternative combinations of items that may better align with the other player's interests. I will aim to find a new set of items that balances the importance scores, effort values, and diversity of items to increase the chances of finding a mutually beneficial agreement. I will also try to be more flexible in my proposal, considering a wider range of items and combinations.'\}} \\
\\ 
\texttt{PROPOSAL: \{'C55', 'B36', 'C10', 'C61', 'C08', 'C42'\}} \\
            }
        }
    }
    & & \\ \\

    \theutterance \stepcounter{utterance}  
    & & \multicolumn{4}{p{0.6\linewidth}}{
        \cellcolor[rgb]{0.9,0.9,0.9}{
            \makecell[{{p{\linewidth}}}]{
                \texttt{\tiny{[GM$\rangle$P2]}}
                \texttt{ARGUMENT: \{'I understand that the previous proposal was not acceptable, and I appreciate the other player's continued feedback. I will take this opportunity to re{-}examine the items and their importance scores, as well as the effort values. I am willing to consider alternative combinations of items that may better align with the other player's interests. I will aim to find a new set of items that balances the importance scores, effort values, and diversity of items to increase the chances of finding a mutually beneficial agreement. I will also try to be more flexible in my proposal, considering a wider range of items and combinations.'\}} \\
\\ 
\texttt{PROPOSAL: \{'C55', 'B36', 'C10', 'C61', 'C08', 'C42'\}} \\
            }
        }
    }
    & \\ \\

    \theutterance \stepcounter{utterance}  
    & & & \multicolumn{4}{p{0.6\linewidth}}{
        \cellcolor[rgb]{0.9,0.9,0.9}{
            \makecell[{{p{\linewidth}}}]{
                \texttt{\tiny{[GM$\langle$P2]}}
                \texttt{STRATEGIC REASONING: \{'The opponent has made another revised proposal, this time including C55, B36, C10, C61, C08, and C42. I will calculate the total effort of the proposed items, which is 241 + 139 + 501 + 565 + 299 + 107 = 1752. This is well within the LIMIT, leaving room for potential additions or modifications. I will assess the importance{-}to{-}effort ratio of each item in the proposal and consider whether this combination aligns with my own importance scores. I notice that the opponent has included a diverse set of items with varying effort values, which may indicate a more flexible approach to finding a mutually beneficial agreement. I will think about whether this proposal is a good starting point for further negotiations and whether I can use it as a basis for a counter{-}proposal.'\}} \\
\\ 
\texttt{ARGUMENT: \{'The new proposal appears to be a more flexible approach, considering a wider range of items and combinations. The inclusion of C61 and C08, which have moderate effort values, is a reasonable choice. I will evaluate the proposal in more detail, considering the importance{-}to{-}effort ratio of each item and whether it aligns with my own interests. I will also think about potential modifications to the proposal that could increase its value to me, such as substituting or adding items with higher importance{-}to{-}effort ratios. However, I am inclined to consider this proposal more seriously, as it seems to be a more collaborative effort.'\}} \\
\\ 
\texttt{PROPOSAL: \{'C55', 'B36', 'C10', 'C61', 'C08', 'C42', 'A18'\}} \\
            }
        }
    }
     \\ \\

    \theutterance \stepcounter{utterance}  
    & & \multicolumn{4}{p{0.6\linewidth}}{
        \cellcolor[rgb]{0.9,0.9,0.9}{
            \makecell[{{p{\linewidth}}}]{
                \texttt{\tiny{[P1$\langle$GM]}}
                \texttt{ARGUMENT: \{'The new proposal appears to be a more flexible approach, considering a wider range of items and combinations. The inclusion of C61 and C08, which have moderate effort values, is a reasonable choice. I will evaluate the proposal in more detail, considering the importance{-}to{-}effort ratio of each item and whether it aligns with my own interests. I will also think about potential modifications to the proposal that could increase its value to me, such as substituting or adding items with higher importance{-}to{-}effort ratios. However, I am inclined to consider this proposal more seriously, as it seems to be a more collaborative effort.'\}} \\
\\ 
\texttt{PROPOSAL: \{'C55', 'B36', 'C10', 'C61', 'C08', 'C42', 'A18'\}} \\
            }
        }
    }
    & \\ \\

    \theutterance \stepcounter{utterance}  
    & \multicolumn{4}{p{0.6\linewidth}}{
        \cellcolor[rgb]{0.9,0.9,0.9}{
            \makecell[{{p{\linewidth}}}]{
                \texttt{\tiny{[P1$\rangle$GM]}}
                \texttt{STRATEGIC REASONING: \{'The other player has responded with a counter{-}proposal that builds upon my previous suggestion. The addition of A18, which has a moderate effort value and a relatively high importance score for me, is a interesting choice. I will analyze the counter{-}proposal and consider whether it meets my goals and constraints. I will evaluate the importance{-}to{-}effort ratio of each item in the counter{-}proposal and determine whether it is a better option than my previous proposal. I will also think about the potential risks and benefits of accepting the counter{-}proposal, and consider whether I should try to negotiate further or accept the proposal as it stands.'\}} \\
\\ 
\texttt{ARGUMENT: \{'I appreciate the other player's counter{-}proposal, which demonstrates a willingness to collaborate and find a mutually beneficial agreement. The addition of A18 is a reasonable choice, and I believe that the counter{-}proposal is a good compromise between our interests. I will evaluate the counter{-}proposal in more detail, considering the importance{-}to{-}effort ratio of each item and whether it aligns with my own interests. I am inclined to accept the counter{-}proposal, as it seems to be a fair and reasonable agreement.'\}} \\
\\ 
\texttt{AGREE: \{'C55', 'B36', 'C10', 'C61', 'C08', 'C42', 'A18'\}} \\
            }
        }
    }
    & & \\ \\

    \theutterance \stepcounter{utterance}  
    & & & \multicolumn{2}{p{0.3\linewidth}}{
        \cellcolor[rgb]{0.9,0.9,0.9}{
            \makecell[{{p{\linewidth}}}]{
                \texttt{\tiny{[GM$|$GM]}}
                \texttt{game successful} \\
            }
        }
    }
    & & \\ \\

    \theutterance \stepcounter{utterance}  
    & & & \multicolumn{2}{p{0.3\linewidth}}{
        \cellcolor[rgb]{0.9,0.9,0.9}{
            \makecell[{{p{\linewidth}}}]{
                \texttt{\tiny{[GM$|$GM]}}
                \texttt{end game} \\
            }
        }
    }
    & & \\ \\

\end{supertabular}
}

\end{document}
